\section{Xây dựng AI Runtime}

\subsection{Điểm khởi chạy}

Tệp \path{apps/edge/edge-runtime/main.py} là entrypoint của container runtime.
\texttt{RuntimeSettings} đọc các biến môi trường do Agent truyền vào, sau đó chọn
adapter: \texttt{source=synthetic} dùng \texttt{SyntheticFrameSource}, các giá trị
khác (URL RTSP, chỉ số camera USB, đường dẫn video) dùng \texttt{OpenCvFrameSource};
có \texttt{MODEL\_PATH} thì dùng \texttt{UltralyticsDetector} với ngưỡng
\texttt{conf} (mặc định 0,25), ngược lại dùng \texttt{SyntheticDetector}. Kết quả
được gửi qua \texttt{HttpTelemetrySink} tới API cục bộ của Agent.

\subsection{Suy luận YOLO}

\texttt{UltralyticsDetector} nạp model một lần khi khởi động bằng
\texttt{YOLO(model\_path)}. Cùng một mã nguồn chạy được checkpoint PyTorch
\texttt{.pt}, mô hình ONNX \texttt{.onnx} hoặc TensorRT engine \texttt{.engine};
định dạng cụ thể được quyết định bởi tệp model trong deployment. Nhờ vậy việc
chuyển từ model FP32 sang engine INT8 sau tối ưu (Chương~\ref{chap:toi-uu}) chỉ
là một lệnh \texttt{update} đổi \texttt{model\_uri}, không cần build lại image.
Mỗi detection gồm chỉ số lớp, tên lớp, độ tin cậy và hộp bao \texttt{xyxy}.

\subsection{Đo lường}

\texttt{InferenceMetrics} (150 dòng) được đặt ở tầng domain để kiểm thử độc lập.
Lớp này ghi thời gian nhận khung, độ sáng, số khung bị bỏ, latency từng lần suy
luận, số lỗi và timeout, các detection; khi \texttt{snapshot()} được gọi, nó tính
FPS theo thời gian thực của cửa sổ, phân vị latency bằng sắp xếp mẫu, thống kê
confidence, entropy của phân bố lớp và tốc độ thay đổi đầu ra (khoảng cách
biến phân toàn phần giữa phân bố lớp của hai cửa sổ liên tiếp), rồi đặt lại bộ
đếm cho cửa sổ kế tiếp. Các quy tắc suy ra trạng thái được hiện thực trực tiếp:

\begin{itemize}
  \item Camera DISCONNECTED khi mất nguồn; FROZEN khi khung không đổi quá 2~s;
        DEGRADED khi FPS camera thấp hơn 70\% FPS mục tiêu; còn lại ONLINE.
  \item Runtime DEGRADED khi tỷ lệ lỗi suy luận vượt 10\% hoặc
        \texttt{effective\_fps\_ratio} dưới 0,5.
\end{itemize}

Mỗi 5~s một mẫu \texttt{RuntimeTelemetry} được gửi tới Agent.

\subsection{Đóng gói}

Image runtime (\path{infra/docker/edge/Dockerfile.worker}) dựng trên
\texttt{python:3.12-slim}, cài thư viện hệ thống cho OpenCV và cài bản Ultralytics
đã được mở rộng trong đồ án (gói \texttt{packages/ultralytics}, chứa các module
pruned và QAT) thay vì bản phát hành trên PyPI, để runtime nạp được checkpoint
đã cắt tỉa. Trên Jetson, base image được thay bằng \texttt{l4t-pytorch} có sẵn
CUDA, cuDNN và TensorRT đúng phiên bản JetPack để engine TensorRT tạo trên thiết
bị chạy được.
